\section{1. Introduction}\label{introduction}

The cryptography deployed on the open Internet today was designed under
assumptions that will not survive the next decade of computing. Shor's
algorithm on a sufficiently large fault-tolerant quantum computer breaks
RSA, finite-field Diffie--Hellman, and every deployed-at-scale elliptic
curve scheme in polynomial time {[}1{]}. Grover's algorithm halves the
effective key length of symmetric primitives {[}2{]}. The exact arrival
date of a \emph{cryptographically relevant quantum computer} (CRQC)
remains contested, but two operational facts are no longer in doubt.
First, adversaries can record encrypted traffic now and decrypt it later
once a CRQC is available --- the \emph{harvest-now, decrypt-later}
threat {[}3{]}, {[}4{]}. Second, the standards bodies have moved: NIST
finalized FIPS 203 (ML-KEM, née Kyber), FIPS 204 (ML-DSA, née
Dilithium), and FIPS 205 (SLH-DSA, née SPHINCS+) in 2024 {[}5{]},
{[}6{]}, {[}7{]}. The U.S. NSA's CNSA 2.0 suite gives calendar
deadlines: software/firmware signing in PQC by 2025, browsers and
servers by 2030, network equipment by 2030, with operating systems and
other classes by 2033 {[}3{]}. The U.K., EU, and ANSSI guidance follows
similar curves {[}8{]}, {[}9{]}.

This places operators in an unfamiliar position. They must inventory,
classify, and migrate cryptographic assets across networks they did not
entirely design, on a deadline they did not set, using algorithms that
did not exist in deployable form three years ago. Tools for the first
step --- inventory and classification --- are appearing {[}10{]},
{[}11{]}, {[}12{]}, {[}13{]}, and standards bodies have begun publishing
guidance on which observables matter {[}4{]}, {[}14{]}.

In parallel, a hardware-rooted track has matured. Quantum Key
Distribution (QKD), proposed in 1984 {[}15{]} and matured through the
SECOQC, Tokyo, and EuroQCI testbeds {[}16{]}, {[}17{]}, {[}18{]},
delivers symmetric key material through a channel whose security reduces
to physics rather than computational assumptions. ETSI's \emph{GS QKD
014} specification standardized the REST interface between QKD
key-management entities (KMEs) and the applications that consume their
keys {[}19{]}. Commercial deployments exist in finance, government, and
metropolitan-fiber settings {[}20{]}, {[}21{]}, {[}22{]}.

Industrial and academic literature treats these two tracks --- PQC and
QKD --- as alternatives. Either you migrate your algorithms (PQC), or
you install QKD links on high-assurance segments, with the implicit
understanding that QKD's hardware cost and distance limits make it
inappropriate for general use {[}23{]}, {[}24{]}. That framing collapses
too soon. For an operator trying to know what risks are present in their
environment on any given day, PQC and QKD are \emph{complementary
observables}, not alternatives. An asset using ECDSA-P256 over a fiber
link protected by a Toshiba MU-QKD3 system is in a materially different
posture than an asset using ECDSA-P256 over an unprotected wide-area
link, even though their algorithms are identical. A unified
observability model must capture both.

A complete quantum-risk model has to include a third axis that current
tooling almost entirely ignores: \emph{crypto-agility}. An asset using a
classical algorithm today but whose algorithm is selected at runtime by
a server configuration field is in a very different posture than an
asset whose algorithm is hard-coded into firmware that ships from a
vendor and rotates only on hardware refresh. The first is a
configuration change; the second is a hardware-replacement program. The
NIST PQ-migration playbook already names crypto-agility as a
precondition for orderly migration {[}4{]}, {[}25{]}, yet no production
observability tool we have seen surfaces it as a first-class telemetry
axis. Without it, operators answering ``are we PQ-ready?'' answer a much
easier question than ``can we \emph{become} PQ-ready in time?''

We argue for, and implement, a \textbf{three-axis posture model}:

\begin{itemize}
\tightlist
\item
  \textbf{Axis A --- Algorithmic Resistance.} Is the primitive itself
  quantum-resistant under standard assumptions?
\item
  \textbf{Axis C --- Channel Protection.} Is the key material protected
  by a quantum-secure delivery mechanism (QKD, hybrid PSK derived from
  QKD)?
\item
  \textbf{Axis G --- Migration Agility.} How quickly can axis A be
  changed on this asset --- by configuration, library upgrade, or
  hardware replacement?
\end{itemize}

We collapse the three into a single deadline-adjusted quantum-risk score
\(q(\mathit{asset}, t)\) suitable for dashboard display, alert
thresholds, and inter-organization comparison. We extend the open Sezar
\texttt{crypto\_inventory\_event\ v1} schema with the additional axes,
implement five agents that emit the same event shape, and evaluate the
system through three empirical studies any practitioner can replicate on
a Linux host with a public Internet connection.

\subsection{1.1 Contributions}\label{contributions}

This paper makes the following contributions:

\begin{enumerate}
\def\labelenumi{\arabic{enumi}.}
\item
  \textbf{A three-axis quantum-risk posture model.} Algorithmic
  resistance, channel protection, and migration agility are three
  independent observables that together determine an asset's actual
  exposure to the post-quantum transition. We formalize each axis on a
  defined scale and define the unified deadline-adjusted risk function
  (§5).
\item
  \textbf{An extended open event schema.} We extend the Sezar
  \texttt{crypto\_inventory\_event\ v1} schema with channel-protection
  and agility fields, and we add new asset kinds for QKD link and KME
  observations. The extensions are additive and non-breaking (§6).
\item
  \textbf{An open observability platform.} We release Sezar --- five
  agents (network, certificate, blockchain, key-management, QKD), a
  posture rollup library, and a collector/dashboard. The implementation
  choices that allow a single Rust workspace to host wire-level eBPF
  observation, REST-based QKD telemetry collection, and static
  crypto-agility analysis under one event shape are detailed in §7.
\item
  \textbf{An ETSI GS QKD 014 emulator and replay corpus.} We provide an
  open implementation of the ETSI GS QKD 014 v1.1.1 Key Delivery API
  with synthetic key generation, configurable QBER, link state-change
  replay scenarios, and a documented capture format. Practitioners
  without QKD hardware can drive and test QKD-aware software against it
  (§7.2, §8.2).
\item
  \textbf{A reproducible crypto-agility scoring rubric.} We define a
  five-level ordinal scale (Negotiated / Configurable / Pinned / Locked
  / Frozen), publish a Semgrep rule pack that derives a defensible score
  from source code or installed-package inspection, and release a
  hand-graded ground-truth corpus of fifty open-source server projects
  (§7.3, §8.3).
\item
  \textbf{Empirical baseline across the three axes.} We report A, C, G
  scores for the Tranco-top-1k over TLS, for a controlled multi-KME ETSI
  014 testbed exercising representative failure modes, and for an
  eleven-project pilot drawn from the fifty-project agility corpus (§8).
  Scaling the pilot to the full corpus is the remaining mechanical step.
\end{enumerate}

\subsection{1.2 Non-goals}\label{non-goals}

This paper does not propose a new PQC primitive, a new QKD protocol, or
a new key-management protocol. It does not argue that QKD should replace
PQC, or vice versa; we take the standards-body position that hybrid
deployments are the realistic operating mode and instrument both
observables. We do not claim to inventory all deployed cryptography on
Earth; our empirical work is bounded by the corpora we publish. We do
not address adversary modeling beyond accepting the
harvest-now/decrypt-later assumption.

\subsection{1.3 Roadmap}\label{roadmap}

§2 reviews the standards landscape (PQC, QKD, crypto-agility) at the
level required to make the paper self-contained. §3 surveys related
observability work and identifies the gap. §4 states the threat model
and operator assumptions. §5 defines the three-axis posture model. §6
specifies the event schema. §7 presents the Sezar reference
implementation. §8 reports the empirical evaluation. §9 discusses
limitations, ethics, and deployment guidance. §10 concludes.

\begin{center}\rule{0.5\linewidth}{0.5pt}\end{center}

\section{2. Background}\label{background}

\subsection{2.1 Post-Quantum Cryptography
Standardization}\label{post-quantum-cryptography-standardization}

NIST initiated its post-quantum cryptography standardization process in
2016 {[}26{]} and finalized the first three production standards in
August 2024 {[}5{]}, {[}6{]}, {[}7{]}. These are:

\begin{itemize}
\tightlist
\item
  \textbf{FIPS 203 --- ML-KEM} (formerly CRYSTALS-Kyber), a
  module-lattice key-encapsulation mechanism at three parameter sets
  corresponding to NIST security categories 1, 3, and 5.
\item
  \textbf{FIPS 204 --- ML-DSA} (formerly CRYSTALS-Dilithium), a
  module-lattice digital signature scheme also at three parameter sets.
\item
  \textbf{FIPS 205 --- SLH-DSA} (formerly SPHINCS+), a stateless
  hash-based digital signature scheme, with both small and fast
  parameter sets across SHA2 and SHAKE hashes.
\end{itemize}

A fourth standard for a fast falcon-style lattice signature (FN-DSA,
formerly Falcon) is in advanced draft {[}27{]}. NIST has additionally
signaled that further KEM candidates will be standardized to diversify
mathematical assumptions {[}28{]}. The U.S. NSA's \emph{Commercial
National Security Algorithm Suite 2.0} (CNSA 2.0) identifies
ML-KEM-1024, ML-DSA-87, SLH-DSA-SHA2-256s, AES-256-GCM, and SHA-384 /
SHA-512 as the required suite for national-security systems by
milestones falling between 2025 and 2035 depending on equipment class
{[}3{]}. The German BSI, French ANSSI, and U.K. NCSC have published
equivalent guidance with broadly similar timelines and a shared
recommendation to deploy \emph{hybrid} (classical + PQC) key exchange
during the transition window {[}8{]}, {[}9{]}, {[}14{]}.

The deployment story is mixed. Chrome and Firefox have shipped
X25519MLKEM768 hybrid key exchange in TLS 1.3 since 2024 {[}29{]},
{[}30{]}. Cloudflare reported roughly two percent of TLS 1.3 connections
to its edge were post-quantum-protected in early 2024 {[}13{]}; by late
2025 the majority of human-initiated traffic to the same edge was
hybrid-PQ {[}31{]}. Linux kernel and OpenSSH support is staged across
recent releases. Certificate authorities have piloted ML-DSA signing in
test hierarchies but have not yet issued production-trust-anchor PQ
certificates {[}32{]}. The pattern is consistent across operator
surveys: the algorithms exist, library support is materializing, but the
long tail of embedded, appliance, and legacy software remains
unaddressed {[}10{]}, {[}11{]}.

\subsection{2.2 Quantum Key Distribution and ETSI GS QKD
014}\label{quantum-key-distribution-and-etsi-gs-qkd-014}

Quantum key distribution, in its prepare-and-measure form, was
introduced by Bennett and Brassard in 1984 {[}15{]} and has been
deployed in operational testbeds for two decades {[}16{]}, {[}17{]}. A
QKD system distributes symmetric key material between two endpoints
using single-photon (or weak-coherent) quantum states, with security
guarantees derived from the no-cloning theorem and the disturbance of
measurement. Modern systems achieve metropolitan ranges (≤100 km dark
fiber) at key rates from kilobits to megabits per second {[}20{]},
{[}21{]}. Trusted-node relaying, satellite QKD {[}33{]}, and twin-field
protocols {[}34{]} extend the reach, though at the cost of additional
assumptions.

QKD's controversial status in the policy community deserves brief
treatment because it informs our observability design choices. The U.S.
NSA has stated that QKD is not a replacement for PQC for
national-security systems, citing implementation-attack surface, limited
authentication coverage, and operational complexity {[}23{]}. The French
ANSSI takes a similar position {[}24{]}. In contrast, the U.K., German,
Italian, and Chinese research and standards communities continue to
invest in QKD infrastructure, including the EU's EuroQCI program
{[}18{]} and multiple national QKD networks. The pragmatic operator
view, which this paper adopts, is that QKD will be deployed on a subset
of high-assurance links (financial inter-datacenter, intelligence,
government inter-site) regardless of which policy community is
``correct,'' and that the deployment will be heterogeneous,
vendor-specific, and operationally opaque without explicit observability
tooling.

ETSI's \emph{GS QKD 014 v1.1.1} specification (2019) standardizes the
REST interface between a Key Management Entity (KME) --- the device that
holds key material produced by the underlying QKD link --- and the
applications consuming those keys (Secure Application Entities, SAEs)
{[}19{]}. The interface defines three operations:

\begin{itemize}
\tightlist
\item
  \texttt{GET\ /api/v1/keys/\{slave\_SAE\_ID\}/status} returns KME and
  link status: current key rate, number of stored keys, supported key
  sizes.
\item
  \texttt{GET\ /api/v1/keys/\{slave\_SAE\_ID\}/enc\_keys} returns one or
  more fresh keys from the master KME, identified by UUIDs.
\item
  \texttt{POST\ /api/v1/keys/\{master\_SAE\_ID\}/dec\_keys} retrieves
  the matching keys at the slave KME, by UUID.
\end{itemize}

The specification is deliberately silent on what an SAE does with the
keys; in practice, SAEs use them as pre-shared keys (PSK) for IPsec,
MACsec, or --- increasingly --- as the symmetric secret in a hybrid TLS
PSK mode {[}35{]}, {[}36{]}. From an observability standpoint, ETSI 014
gives us a stable, vendor-neutral surface on which to observe
QKD-protected channels: link status, key rate, request volume, error
rates, and (with cooperating SAE instrumentation) which downstream
sessions consume the keys.

\subsection{2.3 Crypto-Agility}\label{crypto-agility}

Crypto-agility --- the property that allows a system to change
cryptographic primitives without invasive redesign --- was codified
operationally by RFC 7696 {[}25{]} and is repeatedly named in
PQ-migration guidance as a precondition for orderly transition {[}4{]},
{[}8{]}, {[}14{]}. The literature describes crypto-agility along three
dimensions: \emph{protocol} agility (TLS 1.3 negotiates ciphersuites;
IPsec IKEv2 negotiates transforms), \emph{code} agility (a library
exposes algorithm choice as a parameter, not a recompile-time decision),
and \emph{deployment} agility (operators can roll out a new algorithm
without firmware replacement or trust-anchor reissuance).

In practice, agility varies dramatically across asset classes. A modern
Web server typically negotiates its TLS ciphersuite on every handshake;
a hardware HSM may lock its supported key types at the firmware level;
embedded devices in industrial or medical settings often have a fixed
algorithm chosen at the silicon level. The emergence of FIPS 140-3
validation, with its requirement that algorithm changes trigger
revalidation, introduces a second axis of \emph{locked} status that any
regulated environment has to plan around {[}37{]}.

Despite repeated standards-body emphasis on crypto-agility, the
practitioner picture is poor. No widely deployed inventory tool we have
seen reports crypto-agility as a first-class field. The question ``can
this asset migrate?'' is repeatedly answered informally --- by tribal
knowledge, vendor questionnaires, or case-by-case engineering audits.

\begin{center}\rule{0.5\linewidth}{0.5pt}\end{center}

\section{3. Related Work}\label{related-work}

We organize related work along three threads --- PQC discovery tooling,
QKD telemetry, and crypto-agility analysis --- and conclude with the gap
statement that motivates the present paper.

\subsection{3.1 PQC Discovery and Migration
Tooling}\label{pqc-discovery-and-migration-tooling}

Cisco's PQC Discovery service {[}10{]} is, to our knowledge, the most
operationally mature PQ inventory tool deployed today. It combines
passive network observation (NetFlow-derived context with optional
inline TLS handshake inspection) with active endpoint scanning and a
managed-service operator dashboard. The analytic surface is
algorithm-class oriented: assets are labeled \emph{PQ-ready},
\emph{classical}, or \emph{unknown} based on observed handshakes and
reported algorithm support. The service does not surface
channel-protection or agility as separate dimensions.

IBM's \emph{Quantum-Safe Discover} {[}11{]} and Microsoft's internal
\emph{CryptoTracker} tooling {[}12{]} follow similar patterns: enumerate
cryptographic uses across an environment, classify by algorithm against
a PQC-readiness yardstick, prioritize migration. Cloudflare's published
reports {[}13{]} focus on edge-observable TLS handshake mixtures and
provide one of the few public datasets on PQ adoption on the open
Internet.

The academic literature includes systematic measurements of TLS
deployment hygiene (key sizes, cipher choices, validation behavior)
{[}38{]}, {[}39{]}, {[}40{]} and more recent measurements of hybrid PQ
deployment in the wild {[}31{]}. NIST IR 8547 {[}4{]} provides the most
authoritative migration playbook and includes a discovery section, but
stops short of prescribing telemetry semantics.

Common across this thread: algorithm-class is treated as the primary
axis, hybrid/QKD is at most a footnote, and agility is discussed in
prose but not surfaced as a measured field.

\subsection{3.2 QKD Telemetry}\label{qkd-telemetry}

QKD telemetry has been studied primarily from the perspective of the QKD
operator --- the entity running the physical link and the KMEs. The ETSI
ISG-QKD has produced multiple documents on operational and security
testing {[}41{]}, {[}42{]}, {[}43{]}; the SECOQC {[}16{]} and Tokyo
{[}17{]} testbeds published detailed operational data. The EuroQCI
program {[}18{]} is in the process of standardizing inter-domain QKD
telemetry.

This thread, however, is \emph{internal} to the QKD operator. For an
operator using QKD keys to protect application traffic --- the
SAE/application perspective --- the public literature is much thinner.
ETSI GS QKD 014 {[}19{]} standardizes the SAE-facing interface; we are
not aware of an existing open-source SAE-side observability platform
that consumes 014 telemetry and integrates it with a broader
cryptographic inventory.

\subsection{3.3 Crypto-Agility Analysis}\label{crypto-agility-analysis}

RFC 7696 {[}25{]} is the canonical statement. Subsequent work has
analyzed agility properties of individual protocols (TLS, IPsec, S/MIME)
and proposed agility frameworks for specific domains (industrial
control, embedded), though without operationalising them as a continuous
telemetry input.

In the practitioner space, NCC Group has published audit reports on
cryptographic-agility deficiencies in deployed products {[}44{]}. The
OWASP and CIS communities maintain checklists {[}45{]}, but none of this
work, to our knowledge, has been operationalized as a continuous
telemetry input feeding a cryptographic-posture dashboard.

\subsection{3.4 Gap Statement}\label{gap-statement}

Each thread is mature in isolation. What is missing --- and what Sezar
fills --- is a tool that treats algorithmic resistance, channel
protection, and migration agility as three independent telemetry axes,
rolls them into a single deadline-adjusted score from a shared event
schema, and is reproducible without commercial tooling.

\begin{center}\rule{0.5\linewidth}{0.5pt}\end{center}

\section{4. Threat Model and Operator
Assumptions}\label{threat-model-and-operator-assumptions}

We assume a \emph{Q-day} adversary: a sufficiently capable
fault-tolerant quantum computer capable of executing Shor's algorithm on
RSA-2048-class moduli and elliptic-curve groups of practical size. The
adversary's \emph{capability date} is unknown but treated as a moving
deadline; the standards-body consensus places the planning horizon
between 2030 and 2040 {[}3{]}, {[}4{]}, {[}8{]}. Following the
harvest-now/decrypt-later threat model {[}3{]}, we assume that an
adversary records classical-encrypted traffic \emph{today} and decrypts
it \emph{later} once the CRQC is available. Consequently, traffic
encrypted with classical-only KEMs is at risk \emph{now}, not at Q-day.

We assume an operator who:

\begin{itemize}
\tightlist
\item
  Controls or has visibility into the cryptographic surfaces of their
  environment (TLS termination points, certificate inventory,
  application source code, host configuration, QKD links if any).
\item
  Operates under a calendar deadline (NIST CNSA 2.0 milestones,
  regulatory mandates, or contractual obligations) that is fixed, even
  if individual subdeadlines are uncertain.
\item
  Cannot replace all cryptography at once; migration is staged over
  months to years.
\item
  Tolerates and benefits from continuous telemetry rather than periodic
  audits.
\end{itemize}

We assume an adversary who:

\begin{itemize}
\tightlist
\item
  Records traffic now, decrypts later.
\item
  Cannot break PQ-secure primitives at standard parameter sizes (we
  accept the NIST standardization process's hardness assumptions).
\item
  Cannot break the no-cloning theorem (we accept the standard QKD
  security argument); but may exploit implementation-level QKD
  vulnerabilities --- detector blinding, side channels, classical
  authentication failure --- and so QKD protection is observed
  \emph{probabilistically} by the SAE, not guaranteed.
\item
  May exploit non-cryptographic vulnerabilities; this paper does not
  address those.
\end{itemize}

The threat model justifies treating agility as a \emph{risk multiplier}:
an asset with low algorithmic resistance but high agility may not be at
risk in \emph{real} terms because it can be migrated before the deadline
horizon. An asset with low resistance and low agility (e.g., a hardware
appliance with hard-coded ECDSA on a five-year refresh cycle) is at
\emph{significantly higher} real risk, even if its observed primitive is
identical to the agile case.

\begin{center}\rule{0.5\linewidth}{0.5pt}\end{center}

\section{5. The Three-Axis Quantum-Risk Posture
Model}\label{the-three-axis-quantum-risk-posture-model}

We define three independent axes --- \emph{A}, \emph{C}, \emph{G} ---
and a unified deadline-adjusted quantum-risk score
\(q(\mathit{asset}, t)\).

\subsection{5.1 Axis A --- Algorithmic
Resistance}\label{axis-a-algorithmic-resistance}

Axis \(A\) scores the cryptographic primitives observed on an asset on
the scale \([0, 1]\), where 0 corresponds to a primitive that is
quantum-trivial (a deprecated or known-broken algorithm) and 1
corresponds to a primitive that is quantum-resistant under the standard
NIST hardness assumptions.

Following the existing Sezar V1 rollup (defined in
\texttt{docs/posture-rollup.md} in the Sezar repository), we classify
each primitive into one of five categories:

\begin{longtable}[]{@{}
  >{\raggedright\arraybackslash}p{(\linewidth - 4\tabcolsep) * \real{0.4000}}
  >{\raggedright\arraybackslash}p{(\linewidth - 4\tabcolsep) * \real{0.3143}}
  >{\raggedright\arraybackslash}p{(\linewidth - 4\tabcolsep) * \real{0.2857}}@{}}
\toprule\noalign{}
\begin{minipage}[b]{\linewidth}\raggedright
Category
\end{minipage} & \begin{minipage}[b]{\linewidth}\raggedright
\(a\) value
\end{minipage} & \begin{minipage}[b]{\linewidth}\raggedright
Examples
\end{minipage} \\
\midrule\noalign{}
\endhead
\bottomrule\noalign{}
\endlastfoot
\texttt{pq} & 1.0 & ML-KEM-\{512,768,1024\}, ML-DSA-\{44,65,87\},
SLH-DSA-*, AES-256-GCM, SHA-256 \\
\texttt{pq\_hybrid} & 0.9 & X25519+ML-KEM-768, ECDH-P256+ML-KEM-768 \\
\texttt{classical} & 0.3 & X25519, Ed25519, RSA-2048, ECDSA-P256,
AES-128-GCM \\
\texttt{unknown} & 0.4 & Anything not in the classification table \\
\texttt{deprecated} & 0.0 & SHA-1, MD5, RSA-1024, RC4, 3DES, DH-1024 \\
\end{longtable}

For an asset with observed primitives \(\{p_i\}\) at roles \(\{r_i\}\)
with role weights \(w_{r_i}\) summing to 1, the axis-A score is

\[
A(\mathit{asset}) = \sum_i w_{r_i} \cdot a(p_i)
\]

Role weights are calibrated to harvest-now/decrypt-later risk:
\(w_\text{sig}=0.40, w_\text{kex}=0.30, w_\text{encrypt}=0.20,
w_\text{hash}=0.10\) for assets exhibiting all four roles, with
re-normalization when only a subset is present. The schema also defines
a fifth role, \texttt{auth} --- the MAC primitive in protocols where it
is independent of the AEAD (SSH, IPsec-AH, TLS 1.2 ciphersuites with
separate HMAC). For those assets we fold the \texttt{auth} weight into
\(w_\text{encrypt}\) at observation time; the role exists in the schema
so non-TLS-1.3 protocols can be represented without forcing AEAD
semantics onto them.

\subsection{5.2 Axis C --- Channel
Protection}\label{axis-c-channel-protection}

Axis \(C\) scores the channel through which key material reaches the
endpoint, on the scale \([0, 1]\), where 0 corresponds to a channel with
no quantum-secure key delivery and 1 corresponds to a channel deriving
its session key entirely from QKD-delivered material. Three categorical
states cover the realistic deployment options:

\begin{longtable}[]{@{}
  >{\raggedright\arraybackslash}p{(\linewidth - 4\tabcolsep) * \real{0.4286}}
  >{\raggedright\arraybackslash}p{(\linewidth - 4\tabcolsep) * \real{0.2619}}
  >{\raggedright\arraybackslash}p{(\linewidth - 4\tabcolsep) * \real{0.3095}}@{}}
\toprule\noalign{}
\begin{minipage}[b]{\linewidth}\raggedright
State
\end{minipage} & \begin{minipage}[b]{\linewidth}\raggedright
\(c\) value
\end{minipage} & \begin{minipage}[b]{\linewidth}\raggedright
Description
\end{minipage} \\
\midrule\noalign{}
\endhead
\bottomrule\noalign{}
\endlastfoot
\texttt{classical} & 0.0 & Session key derived solely from the
negotiated KEM. \\
\texttt{qkd\_hybrid\_psk} & 0.7 & Session key derived from QKD-PSK XOR
negotiated KEM (NIST SP 1800-38A pattern, ETSI 014 SAE). \\
\texttt{qkd\_only} & 1.0 & Session key derived from QKD material alone
(rare; MACsec-style transport). \\
\end{longtable}

Sub-states allow partial credit when telemetry indicates a QKD link is
\emph{degraded} --- high QBER, sustained KME unavailability, low key
rate --- but the SAE has not failed over to classical. We discount the
score in proportion to the observed degradation; formal degradation
thresholds are deferred to the implementation section (§7.2).

A subtle but important observation: the SAE may \emph{think} it is
operating in QKD-hybrid mode while the underlying KME is failing
gracefully to a classical fallback. Sezar's role is to observe both
layers and surface the discrepancy. We define \(c\) on observed ETSI 014
status, not on SAE-reported intent.

\subsection{5.3 Axis G --- Migration
Agility}\label{axis-g-migration-agility}

Axis \(G\) scores the asset's ability to migrate its primitives on the
scale \([0, 1]\), where 0 corresponds to an asset whose primitives can
be changed only by physical replacement and 1 corresponds to an asset
that negotiates primitives on every session. We define five ordinal
levels:

\begin{longtable}[]{@{}
  >{\raggedright\arraybackslash}p{(\linewidth - 4\tabcolsep) * \real{0.2540}}
  >{\raggedright\arraybackslash}p{(\linewidth - 4\tabcolsep) * \real{0.1746}}
  >{\raggedright\arraybackslash}p{(\linewidth - 4\tabcolsep) * \real{0.5714}}@{}}
\toprule\noalign{}
\begin{minipage}[b]{\linewidth}\raggedright
Level
\end{minipage} & \begin{minipage}[b]{\linewidth}\raggedright
\(g\) value
\end{minipage} & \begin{minipage}[b]{\linewidth}\raggedright
Definition (observable signature)
\end{minipage} \\
\midrule\noalign{}
\endhead
\bottomrule\noalign{}
\endlastfoot
\texttt{negotiated} & 1.0 & Algorithm selected per-session by protocol
negotiation. (TLS 1.3 server, modern SSH server, IKEv2 responder.) \\
\texttt{configurable} & 0.75 & Algorithm fixed per-deployment but
changeable by configuration without code change. (Library config file,
environment variable.) \\
\texttt{pinned} & 0.50 & Algorithm fixed in code; changeable by software
upgrade. (Hard-coded algorithm name in application source.) \\
\texttt{locked} & 0.20 & Algorithm fixed in firmware or by
FIPS/compliance binding; changeable only by vendor update or
revalidation cycle. (Embedded firmware crypto, FIPS 140-3 validated
module under tested-configuration constraint.) \\
\texttt{frozen} & 0.0 & Algorithm fixed in silicon, ROM, or otherwise
unchangeable without hardware replacement. (TPM 1.2 with hard-coded
RSA-2048; smart-card hard-wired ECDSA-P256.) \\
\end{longtable}

The classification is derived from static analysis of the asset's
implementation surface plus, where available, vendor declarations of
FIPS validation scope. We detail the scanning methodology in §7.3 and
the scoring rubric in §8.3.

\subsection{5.4 Unified Deadline-Adjusted Quantum
Risk}\label{unified-deadline-adjusted-quantum-risk}

The three axes are independent observables. To collapse them into a
single posture metric we adopt the operator's deadline as a fourth
input.

Let \(D\) be the operator-configured deadline (e.g., 2030-01-01 for NSA
CNSA 2.0 browser/server class). Let \(t\) be the current date. Define
the \emph{deadline tension} \(\tau(t) = \max(0,
\min(1, 1 - (D-t)/H))\), where \(H\) is a horizon constant (we use five
years by default). When \(D\) is far in the future, \(\tau \to 0\) and
agility forgives lower algorithmic resistance. As \(t \to D\),
\(\tau \to 1\) and agility no longer compensates because there is
insufficient time to migrate.

We define the quantum-risk score as

\[
q(\mathit{asset}, t) = 1 - \Bigl( \alpha \cdot A + \beta \cdot C +
  \gamma(\tau) \cdot G \Bigr)
\]

with \(\alpha + \beta + \gamma(\tau) = 1\). Default weights:
\(\alpha = 0.5, \beta = 0.2, \gamma(\tau) = 0.3 \cdot (1-\tau)\),
re-normalized when \(\gamma\) shrinks. The agility weight is the only
weight that shrinks with deadline tension: an asset with high agility
but classical algorithms looks safe today (because \(\gamma\) is large)
and looks increasingly unsafe as the deadline approaches (because
\(\gamma\) shrinks toward zero). This is the intended behavior: agility
is forgiving \emph{now}, not \emph{at the deadline}.

Asset-class weights \(w_k\) further weight the asset's contribution to
the org-wide posture, with \texttt{blockchain\_key} weighted higher than
\texttt{tls\_session} (a forged signature against a public-chain key is
permanent; a forged session is ephemeral).

\subsection{5.5 Interpretation and
Bounds}\label{interpretation-and-bounds}

The score \(q \in [0, 1]\), with 0 = posture is fully aligned with the
deadline and 1 = posture is maximally exposed. Operators typically
configure alert thresholds at \(q > 0.6\) (``must migrate'') and
\(q > 0.3\) (``plan migration'').

Two edge cases warrant note. First, a fully PQ asset on a
non-QKD-protected channel scores \(q = 1 - (0.5 \cdot 1 + 0.2 \cdot
0 + 0.3 \cdot g)\), which is approximately 0.2 for a fully agile asset
and 0.5 for a fully frozen asset, both before deadline tension. This is
intentional: QKD is a \emph{bonus} on PQ-protected sessions, not a
requirement. Second, a fully classical asset on a QKD-protected channel
scores approximately 0.35 even with no agility --- QKD partially
compensates for classical algorithms, reflecting the real-world
deployment of QKD on high-assurance links.

\begin{center}\rule{0.5\linewidth}{0.5pt}\end{center}

\section{6. Event Schema Extensions}\label{event-schema-extensions}

We extend the Sezar \texttt{crypto\_inventory\_event\ v1} schema
(defined in \texttt{docs/crypto-event-schema.md} in the Sezar
repository) with additive, non-breaking fields. Existing consumers that
ignore unknown top-level fields continue to function; consumers that opt
in to v1.1 gain access to channel-protection and agility observables.

\subsection{6.1 New Top-Level Fields}\label{new-top-level-fields}

\begin{Shaded}
\begin{Highlighting}[]
\FunctionTok{\{}
  \DataTypeTok{"schema\_version"}\FunctionTok{:} \DecValTok{1}\FunctionTok{,}
  \DataTypeTok{"schema\_minor"}\FunctionTok{:} \DecValTok{1}\FunctionTok{,}
  \DataTypeTok{"source\_module"}\FunctionTok{:} \StringTok{"sezar{-}net"}\FunctionTok{,}
  \DataTypeTok{"observed\_at"}\FunctionTok{:} \StringTok{"2026{-}08{-}15T11:42:03.421Z"}\FunctionTok{,}
  \DataTypeTok{"asset"}\FunctionTok{:} \FunctionTok{\{} \ErrorTok{...} \FunctionTok{\},}
  \DataTypeTok{"primitives"}\FunctionTok{:} \OtherTok{[} \ErrorTok{...} \OtherTok{]}\FunctionTok{,}
  \DataTypeTok{"channel\_protection"}\FunctionTok{:} \FunctionTok{\{} \ErrorTok{...} \FunctionTok{\},}   \ErrorTok{//} \ErrorTok{NEW} \ErrorTok{in} \ErrorTok{1.1}
  \DataTypeTok{"agility"}\FunctionTok{:} \FunctionTok{\{} \ErrorTok{...} \FunctionTok{\},}              \ErrorTok{//} \ErrorTok{NEW} \ErrorTok{in} \ErrorTok{1.1}
  \DataTypeTok{"posture"}\FunctionTok{:} \FunctionTok{\{} \ErrorTok{...} \FunctionTok{\}}
\FunctionTok{\}}
\end{Highlighting}
\end{Shaded}

\texttt{schema\_minor} is the first additive use of a minor-version
field. Consumers must accept any \texttt{schema\_minor} ≥ their compiled
value and treat unknown top-level fields as opaque.

\subsection{\texorpdfstring{6.2
\texttt{channel\_protection}}{6.2 channel\_protection}}\label{channel_protection}

\begin{Shaded}
\begin{Highlighting}[]
\FunctionTok{\{}
  \DataTypeTok{"state"}\FunctionTok{:} \StringTok{"qkd\_hybrid\_psk"}\FunctionTok{,}
  \DataTypeTok{"kme\_endpoint"}\FunctionTok{:} \StringTok{"https://kme{-}1.dc.example/api/v1"}\FunctionTok{,}
  \DataTypeTok{"key\_id\_observed"}\FunctionTok{:} \StringTok{"9c45e0a2{-}..."}\FunctionTok{,}
  \DataTypeTok{"psk\_age\_seconds"}\FunctionTok{:} \DecValTok{47}\FunctionTok{,}
  \DataTypeTok{"link\_qber"}\FunctionTok{:} \FloatTok{0.018}\FunctionTok{,}
  \DataTypeTok{"link\_key\_rate\_bps"}\FunctionTok{:} \DecValTok{12480}\FunctionTok{,}
  \DataTypeTok{"link\_health"}\FunctionTok{:} \StringTok{"ok"}\FunctionTok{,}
  \DataTypeTok{"degraded\_reason"}\FunctionTok{:} \KeywordTok{null}
\FunctionTok{\}}
\end{Highlighting}
\end{Shaded}

\begin{longtable}[]{@{}
  >{\raggedright\arraybackslash}p{(\linewidth - 4\tabcolsep) * \real{0.5946}}
  >{\raggedright\arraybackslash}p{(\linewidth - 4\tabcolsep) * \real{0.2162}}
  >{\raggedright\arraybackslash}p{(\linewidth - 4\tabcolsep) * \real{0.1892}}@{}}
\toprule\noalign{}
\begin{minipage}[b]{\linewidth}\raggedright
Field
\end{minipage} & \begin{minipage}[b]{\linewidth}\raggedright
Type
\end{minipage} & \begin{minipage}[b]{\linewidth}\raggedright
Notes
\end{minipage} \\
\midrule\noalign{}
\endhead
\bottomrule\noalign{}
\endlastfoot
\texttt{state} & enum & \texttt{classical} / \texttt{qkd\_hybrid\_psk} /
\texttt{qkd\_only}. \\
\texttt{kme\_endpoint} & string & ETSI 014 base URL. Omitted when state
= \texttt{classical}. \\
\texttt{key\_id\_observed} & string & UUID of the consumed key, when
reported by the cooperating SAE. \\
\texttt{psk\_age\_seconds} & int & Age of the PSK when the session
began. \\
\texttt{link\_qber} & float & Quantum bit error rate (0--1). \\
\texttt{link\_key\_rate\_bps} & int & Average key generation rate over
the prior minute. \\
\texttt{link\_health} & enum & \texttt{ok} / \texttt{degraded} /
\texttt{failed}. \\
\texttt{degraded\_reason} & string & One-sentence reason when degraded;
\texttt{null} otherwise. \\
\end{longtable}

The fields are populated by \texttt{sezar-qkd}, which polls the ETSI 014
\texttt{/status} endpoint at configurable intervals. SAE-side fields
(\texttt{key\_id\_observed}, \texttt{psk\_age\_seconds}) require
cooperating instrumentation in the SAE; when absent they are
\texttt{null}.

\subsection{\texorpdfstring{6.3
\texttt{agility}}{6.3 agility}}\label{agility}

\begin{Shaded}
\begin{Highlighting}[]
\FunctionTok{\{}
  \DataTypeTok{"level"}\FunctionTok{:} \StringTok{"configurable"}\FunctionTok{,}
  \DataTypeTok{"level\_score"}\FunctionTok{:} \FloatTok{0.75}\FunctionTok{,}
  \DataTypeTok{"evidence"}\FunctionTok{:} \OtherTok{[}
    \FunctionTok{\{}
      \DataTypeTok{"type"}\FunctionTok{:} \StringTok{"config\_pattern"}\FunctionTok{,}
      \DataTypeTok{"file"}\FunctionTok{:} \StringTok{"/etc/nginx/nginx.conf"}\FunctionTok{,}
      \DataTypeTok{"line"}\FunctionTok{:} \DecValTok{142}\FunctionTok{,}
      \DataTypeTok{"snippet"}\FunctionTok{:} \StringTok{"ssl\_protocols TLSv1.2 TLSv1.3;}\CharTok{\textbackslash{}n}\StringTok{ssl\_ciphers HIGH:..."}
    \FunctionTok{\}}\OtherTok{,}
    \FunctionTok{\{}
      \DataTypeTok{"type"}\FunctionTok{:} \StringTok{"fips\_mode"}\FunctionTok{,}
      \DataTypeTok{"detected"}\FunctionTok{:} \KeywordTok{false}
    \FunctionTok{\}}
  \OtherTok{]}\FunctionTok{,}
  \DataTypeTok{"scanner\_version"}\FunctionTok{:} \StringTok{"sezar{-}agility/0.3.1"}\FunctionTok{,}
  \DataTypeTok{"rubric\_version"}\FunctionTok{:} \StringTok{"qra{-}rubric/v1.0"}
\FunctionTok{\}}
\end{Highlighting}
\end{Shaded}

\begin{longtable}[]{@{}
  >{\raggedright\arraybackslash}p{(\linewidth - 4\tabcolsep) * \real{0.5588}}
  >{\raggedright\arraybackslash}p{(\linewidth - 4\tabcolsep) * \real{0.2353}}
  >{\raggedright\arraybackslash}p{(\linewidth - 4\tabcolsep) * \real{0.2059}}@{}}
\toprule\noalign{}
\begin{minipage}[b]{\linewidth}\raggedright
Field
\end{minipage} & \begin{minipage}[b]{\linewidth}\raggedright
Type
\end{minipage} & \begin{minipage}[b]{\linewidth}\raggedright
Notes
\end{minipage} \\
\midrule\noalign{}
\endhead
\bottomrule\noalign{}
\endlastfoot
\texttt{level} & enum & \texttt{negotiated} / \texttt{configurable} /
\texttt{pinned} / \texttt{locked} / \texttt{frozen}. \\
\texttt{level\_score} & float & Numeric value per §5.3. \\
\texttt{evidence} & array & One or more evidentiary findings supporting
the level. \\
\texttt{scanner\_version} & string & Sezar-agility version that produced
the score. \\
\texttt{rubric\_version} & string & Version of the public scoring rubric
(§8.3). \\
\end{longtable}

Evidence types in V1: \texttt{protocol\_negotiation} (observed
wire-level algorithm negotiation), \texttt{config\_pattern}
(configuration file exposing algorithm choice), \texttt{code\_pattern}
(source code reference to a fixed algorithm), \texttt{firmware\_string}
(binary-extracted algorithm name), \texttt{fips\_mode} (FIPS
provider/kernel mode detected), \texttt{vendor\_declaration}
(operator-provided vendor statement).

\subsection{6.4 New Asset Kinds}\label{new-asset-kinds}

\begin{verbatim}
qkd_link          // identity = KME endpoint URL hash
qkd_kme           // identity = KME ID per ETSI 014 status
\end{verbatim}

These are emitted by \texttt{sezar-qkd} independently of the session
events that consume their keys. They allow the dashboard to render a QKD
link health view distinct from the SAE-side session view.

\subsection{6.5 Backwards Compatibility}\label{backwards-compatibility}

The schema extension is strictly additive:

\begin{itemize}
\tightlist
\item
  All fields new in v1.1 are top-level; no existing field shape changes.
\item
  v1.0 consumers that do not recognize the new fields ignore them.
\item
  v1.1 producers that lack data for the new fields emit them as
  \texttt{null} (per the V1 module emission contract).
\item
  The posture engine treats \texttt{channel\_protection:\ null} as
  \texttt{state:\ classical} and \texttt{agility:\ null} as the
  \texttt{UNKNOWN\_LEVEL\_FALLBACK} constant defined in
  \texttt{sezar-agility}, which evaluates to \texttt{level:\ pinned}
  (\texttt{level\_score:\ 0.50}). This is deliberately more conservative
  than Axis A's \texttt{unknown\ =\ 0.4}: an asset whose agility we
  cannot evidence is treated as if it were merely pinned rather than
  freely configurable, so the rollup does not credit unmeasured agility
  for free.
\end{itemize}

\begin{center}\rule{0.5\linewidth}{0.5pt}\end{center}

\section{7. Sezar: Reference Architecture and
Implementation}\label{sezar-reference-architecture-and-implementation}

Sezar is an open-source observability platform implementing the
three-axis posture model. The implementation is a single Rust workspace
containing seven crates: five agents emitting events, one shared rollup
library, and one collector/server.

\subsection{7.1 Workspace Layout}\label{workspace-layout}

\begin{verbatim}
sezar/
├── crates/
│   ├── sezar-core/       # event schema, rollup engine (no I/O)
│   ├── sezar-server/     # axum collector + REST API
│   ├── sezar-net/        # eBPF agent: TLS, SSH, IPsec
│   ├── sezar-qkd/        # ETSI GS QKD 014 collector + emulator
│   ├── sezar-cert/       # X.509 inventory (CT, host scan)
│   ├── sezar-chain/      # public-chain crypto observation
│   ├── sezar-id/         # HSM/KMS/smart-card inventory
│   └── sezar-agility/    # static crypto-agility scanner
├── docs/
└── web/                  # React + Vite dashboard
\end{verbatim}

The architectural invariant --- every agent emits one and only one event
shape, computed locally via \texttt{sezar-core::rollup} --- is what
keeps the platform composable across surfaces as different as eBPF TLS
sniffing and Solidity source code analysis.

\subsection{7.2 sezar-qkd: ETSI 014 Collector and
Emulator}\label{sezar-qkd-etsi-014-collector-and-emulator}

The \texttt{sezar-qkd} crate fulfils two roles. Operationally, it is a
collector that polls one or more ETSI GS QKD 014 KMEs, emits
\texttt{qkd\_link} and \texttt{qkd\_kme} events, and serves as the data
source for the \texttt{channel\_protection} block on session events
emitted by cooperating SAEs.

\subsubsection{7.2.1 Collector design}\label{collector-design}

The collector is a Tokio async loop that, for each configured KME,
issues:

\begin{itemize}
\tightlist
\item
  \texttt{GET\ /api/v1/keys/\{slave\_SAE\_ID\}/status} at a configurable
  cadence (default 5 s).
\item
  An auxiliary \texttt{enc\_keys} request at a slower cadence (default
  60

  \begin{enumerate}
  \def\labelenumi{\alph{enumi})}
  \setcounter{enumi}{18}
  \tightlist
  \item
    to measure end-to-end key delivery latency. Keys are requested with
    \texttt{size=0} when permitted, or discarded when not.
  \end{enumerate}
\end{itemize}

The collector tracks per-KME state (last good status time,
exponentially-weighted error rate, QBER history) and emits
\texttt{qkd\_link} events on status change, plus heartbeat events at a
configurable interval. Authentication to the KME follows ETSI 014
guidance: mutual TLS with the SAE certificate.

\subsubsection{7.2.2 Emulator}\label{emulator}

Hardware QKD is expensive and rare, so we ship
\texttt{sezar-qkd-kme-emulator} --- an ETSI 014 v1.1.1 implementation
backed by a synthetic key generator. The emulator:

\begin{itemize}
\tightlist
\item
  Implements \texttt{/status}, \texttt{/enc\_keys}, and
  \texttt{/dec\_keys} exactly per spec.
\item
  Generates synthetic keys at a configured rate, with configurable QBER,
  key size, and lifetime.
\item
  Supports replay scenarios --- pre-recorded sequences of link state
  changes (degradation, failure, recovery) --- for reproducibility.
\item
  Logs every interaction in a documented JSON capture format enabling
  head-to-head A/B testing of SAE implementations.
\end{itemize}

The emulator is the foundation of the §8.2 empirical study and is
released alongside Sezar as a standalone tool.

\subsection{7.3 sezar-agility: Static Crypto-Agility
Scanner}\label{sezar-agility-static-crypto-agility-scanner}

The \texttt{sezar-agility} crate implements the agility-axis scoring per
§5.3. It accepts as input one or more \emph{targets} and produces
\texttt{agility} blocks attached to the corresponding assets.

\subsubsection{7.3.1 Target types}\label{target-types}

\begin{longtable}[]{@{}
  >{\raggedright\arraybackslash}p{(\linewidth - 2\tabcolsep) * \real{0.5135}}
  >{\raggedright\arraybackslash}p{(\linewidth - 2\tabcolsep) * \real{0.4865}}@{}}
\toprule\noalign{}
\begin{minipage}[b]{\linewidth}\raggedright
Target
\end{minipage} & \begin{minipage}[b]{\linewidth}\raggedright
Evidence sources
\end{minipage} \\
\midrule\noalign{}
\endhead
\bottomrule\noalign{}
\endlastfoot
Source repository & Semgrep ruleset over the language-specific config
and source files; file-path heuristics for build-time pinning. \\
Installed package & Package manifest (rpm, dpkg, pip, npm) + binary
string-extraction over the installed artifacts. \\
Running host & TLS handshake observation against the host (server
algorithm support → \texttt{negotiated} evidence); plus optional auth
into the host's config files. \\
Vendor appliance & Vendor-declared algorithm scope + observed handshake
behavior. \\
\end{longtable}

\subsubsection{7.3.2 Ruleset}\label{ruleset}

The published ruleset (\texttt{sezar-agility/rules/v1}) consists of
several hundred Semgrep patterns covering common cryptographic libraries
and protocols across C, Go, Rust, Python, Java, and configuration
formats (nginx, Apache, OpenSSL config, sshd\_config, strongswan.conf,
Postfix, Dovecot, HAProxy, Envoy). Each pattern emits one piece of
evidence with a documented mapping to the five-level rubric.

\subsubsection{7.3.3 Scoring algorithm}\label{scoring-algorithm}

For each target, the scanner collects evidence and applies a
\textbf{most-agile-wins} aggregation: the asset is scored at the
\emph{most} agile level supported by any rule that fired. The rationale
is asymmetric: rules fall into two semantic classes. \emph{Capability}
rules (e.g., presence of an \texttt{ssl\_ciphers} directive, a
\texttt{SSL\_CTX\_set\_cipher\_list} call) demonstrate that an operator
can in fact change the algorithm without invasive surgery; their
emit\_level is an \emph{upper bound on what is observable}.
\emph{Constraint} rules (e.g., a literal hard-coded algorithm name in
source) only report that \emph{some} algorithm is fixed somewhere ---
they do not imply that there is no overriding config surface.
Conservative-min aggregation systematically misclassifies large agility
projects like nginx (which simultaneously embeds default algorithm
strings in source and exposes ten configuration knobs that override
them); the operationally useful question is ``what is the best surface I
have,'' not ``what is the worst constraint anywhere.''

Conservative-min remains available as an alternative aggregation
projection for operators who require it (regulatory contexts where the
weakest surface governs); the published scanner exposes both projections
through a dashboard toggle.

Where evidence is absent, the scanner produces \texttt{level:\ pinned}
(the documented \texttt{UNKNOWN\_LEVEL\_FALLBACK} constant in
\texttt{sezar-agility}) and surfaces the issue for operator review.

We discuss limitations of static-only agility scoring in §9.

\subsection{7.4 sezar-core Unified
Rollup}\label{sezar-core-unified-rollup}

The rollup engine extends the V1 implementation (see
\texttt{docs/posture-rollup.md} in the Sezar repository) with the
deadline-adjusted three-axis formula of §5.4. The implementation remains
a pure function with no I/O; the operator configures \(D\), \(H\), and
asset-class weights via the dashboard, and the engine recomputes on
every event ingest. Fuzz testing covers all axis combinations.

\subsection{7.5 sezar-server and
Dashboard}\label{sezar-server-and-dashboard}

The collector is an Axum HTTP service accepting v1.0 and v1.1 events,
validating against a generated JSON schema, persisting to Postgres
(configuration) and a columnar store (events), and serving a React+Vite
dashboard. The dashboard renders a three-axis posture matrix per asset
class, a deadline-countdown view tied to the configured \(D\), and an
inventory table sortable by \(q\). Implementation details are routine
and we defer them to the project documentation.

\begin{center}\rule{0.5\linewidth}{0.5pt}\end{center}

\section{8. Empirical Evaluation}\label{empirical-evaluation}

We evaluate the model and implementation through three studies, each
designed to be reproducible by an external practitioner with the
published rulesets, emulator, and corpus lists.

\subsection{8.1 Study 1 --- Axis A on the Public Web
(Tranco-top-1k)}\label{study-1-axis-a-on-the-public-web-tranco-top-1k}

\subsubsection{8.1.1 Methodology}\label{methodology}

We scan the Tranco-top-1k {[}46{]} over TLS with two purpose-built
probes --- general-purpose scanners such as \texttt{zgrab2} {[}47{]} do
not yet advertise X25519MLKEM768 in the ClientHello, so we built
dedicated baseline and PQ-capable probes to keep the ClientHello surface
deterministic and the output schema aligned with Sezar's collector:

\begin{enumerate}
\def\labelenumi{\arabic{enumi}.}
\tightlist
\item
  \textbf{Classical baseline probe} --- Python \texttt{ssl} with the
  system OpenSSL defaults. Establishes a TLS handshake without
  advertising \texttt{X25519MLKEM768}. Output: JSON array.
\item
  \textbf{PQ-capable probe} --- a Rust binary built on
  \texttt{rustls\ =\ 0.23} + \texttt{rustls-post-quantum\ =\ 0.2}, with
  a custom certificate verifier that accepts every chain (this is
  observability, not authentication). Advertises \texttt{X25519MLKEM768}
  alongside the classical groups. Output: one NDJSON record per host.
\end{enumerate}

For each host we record: connect/handshake outcome, negotiated TLS
version, negotiated cipher suite, negotiated key-exchange group, leaf
certificate signature algorithm, and the leaf certificate Subject. The
scan is constrained to one TCP connection per host, identifies itself in
the ClientHello SNI extension as
\texttt{sezar-survey/1.0\ +https://e2esolutions.tech/sezar}, and uses a
5-second connect+handshake timeout with a 1 Hz rate cap.

The scan source code, target list, raw probe outputs, and analysis
script are released alongside this paper at
\href{studies/study1/}{\texttt{studies/study1/}}.

Ethical considerations: this is a one-shot benign TLS handshake similar
to common research scans {[}38{]}; we do not attempt protocol downgrade
or repeated probing. Scan rate is ≤1 Hz per probe and exits immediately
on connection error.

\subsubsection{8.1.2 Metrics}\label{metrics}

We report the distribution of asset-A scores across the corpus, the
prevalence of PQ-capable hosts, the distribution of certificate
signature algorithms, and the prevalence of weak/ deprecated primitives
still observed.

\subsubsection{8.1.3 Results (Tranco-1k snapshot 6G8PX,
2026-05-13)}\label{results-tranco-1k-snapshot-6g8px-2026-05-13}

We ran the methodology of §8.1.1 against the Tranco-top-1k list
\texttt{6G8PX} (snapshot 2026-05-13). 724 of 1,000 hosts returned a
usable TLS handshake within the 5-second timeout; the remaining 276 were
unresponsive (DNS failure, no TCP, no TLS on 443, regional GeoIP block,
or anti-bot middlebox). The 27.6\% non-response rate is in the range
reported by prior large-scale TLS scans of the open Web {[}38{]}. All
headline percentages in this section use the \(n = 724\) responsive
denominator unless stated otherwise.

For sample-selection contrast we also report results from a 30-host
curated pilot --- major CDN, browser-vendor, distro,
IETF/IEEE/NIST/ETSI, and AI-vendor properties --- that ran the same
probe pair before the Tranco-1k scan.

\begin{longtable}[]{@{}
  >{\raggedright\arraybackslash}p{(\linewidth - 4\tabcolsep) * \real{0.4568}}
  >{\raggedleft\arraybackslash}p{(\linewidth - 4\tabcolsep) * \real{0.2716}}
  >{\raggedleft\arraybackslash}p{(\linewidth - 4\tabcolsep) * \real{0.2716}}@{}}
\toprule\noalign{}
\begin{minipage}[b]{\linewidth}\raggedright
Observable
\end{minipage} & \begin{minipage}[b]{\linewidth}\raggedleft
Tranco-1k (n=724)
\end{minipage} & \begin{minipage}[b]{\linewidth}\raggedleft
Curated pilot (n=30)
\end{minipage} \\
\midrule\noalign{}
\endhead
\bottomrule\noalign{}
\endlastfoot
TLS 1.3 negotiation & 602/724 (83.1\%) & 30/30 (100\%) \\
TLS 1.2 fallback & 122/724 (16.9\%) & 0/30 (0\%) \\
AES-256-GCM/SHA-384 (TLS 1.3) & 384/724 (53.0\%) & 20/30 (67\%) \\
AES-128-GCM/SHA-256 (TLS 1.3) & 207/724 (28.6\%) & 7/30 (23\%) \\
ChaCha20-Poly1305 (TLS 1.3) & 11/724 (1.5\%) & 3/30 (10\%) \\
ECDSA leaf cert (P256 + P384) & 254/724 (35.1\%) & 18/30 (60\%) \\
RSA-PKCS1 leaf cert (SHA256+SHA384) & 470/724 (64.9\%) & 12/30 (40\%) \\
ML-DSA / SLH-DSA cert & 0/724 & 0/30 \\
Deprecated primitive (SHA-1, RC4) & 0/724 & 0/30 \\
\textbf{\texttt{X25519MLKEM768} negotiated} & \textbf{317/724 (43.8\%)}
& \textbf{17/30 (57\%)} \\
\end{longtable}

\textbf{The headline PQ-adoption result is 317/724 (43.8\%) on the
Tranco-top-1k.} This is close to the 39\% of top-100k sites that
Cloudflare measured supporting post-quantum key agreement in September
2025 {[}31{]}; both numbers are site-capability measurements, and the
4-point gap is consistent with Tranco-top-1k over-representing large
CDN- and cloud-fronted properties relative to the broader top-100k
corpus. The remaining 407 responsive hosts fell back to classical
\texttt{x25519} (312), \texttt{secp256r1} (80), or \texttt{secp384r1}
(15).

The curated 30-host pilot returned 17/30 = 57\% PQ adoption --- 13
percentage points above the Tranco-1k rate. The gap is attributable to
sample-selection bias: the curated list emphasises Cloudflare-fronted
and Google-fronted properties (cloudflare.com, twitter.com, reddit.com,
anthropic.com, openai.com, google.com, youtube.com, e2esolutions.tech,
and a long tail) whose edge terminators rolled out X25519MLKEM768 early.
The 13 classical-only hosts in the curated pilot --- github.com,
microsoft.com, amazon.com, mozilla.org, kernel.org, nist.gov,
openssl.org, and others --- are notable in their own right as major
infrastructure-relevant properties where the PQ rollout has not yet
landed. The two numbers together suggest that surveys built on curated
``interesting'' host lists overstate PQ readiness compared to broader
corpora. Where these numbers appear in the literature, both rates ---
broad-corpus and curated --- deserve to be reported.

\textbf{Two further observables of operational interest.} First, 122 of
724 responsive Tranco-1k hosts (16.9\%) still negotiate TLS 1.2 at the
top of the Web. Because \texttt{X25519MLKEM768} is a TLS 1.3 key-share,
that 16.9\% cohort is structurally PQ-ineligible until the operator's
TLS terminator can negotiate TLS 1.3 --- a long-tail constraint not
visible in any single-axis cipher dashboard. Second, the
certificate-signature distribution on the Tranco-1k is RSA-dominant
(64.9\%), inverting the ECDSA-dominant (60\%) distribution in the
curated pilot. The curated list skews toward post-2020 CDN deployments
that default to ECDSA; the broader corpus retains a substantial RSA
installed base.

Pipeline integration: PQ-probe NDJSON parses straight into the same
\texttt{tls\_session} asset shape consumed by \texttt{from-zgrab}, so
the same downstream rollup, BLOCKED-flag derivation, and dashboard
render exactly as for any other source-module input.

\subsection{8.2 Study 2 --- Axis C via the ETSI GS QKD 014
Emulator}\label{study-2-axis-c-via-the-etsi-gs-qkd-014-emulator}

\subsubsection{8.2.1 Methodology}\label{methodology-1}

We construct a controlled testbed of one master KME, two slave KMEs, and
three cooperating SAEs (a strongSwan IPsec endpoint operating in PSK
mode with QKD-PSK rotation, a Wireguard endpoint with manual PSK
rotation, and a custom TLS endpoint using the NIST SP 1800-38A
hybrid-PSK pattern). The KMEs are \texttt{sezar-qkd-kme-emulator}
instances; we drive them through a documented sequence of replay
scenarios:

\begin{itemize}
\tightlist
\item
  \textbf{R1 --- Steady-state.} Constant QBER 1.8\%, key rate 12 kbps,
  no interruption. 24 hours.
\item
  \textbf{R2 --- Gradual degradation.} QBER ramps from 1.8\% to 8.5\%
  over 4 hours; the SAE policy should fail over to classical at the
  configured threshold. We observe whether each SAE detects the
  degradation and whether Sezar's \texttt{link\_health} reflects the
  transition.
\item
  \textbf{R3 --- Hard failure.} KME unreachable for 30 minutes. SAE
  behavior should be: continue with cached PSK until lifetime expires,
  then fail closed or fall back to classical per policy.
\item
  \textbf{R4 --- Stale PSK.} PSK rotation is suppressed at the SAE while
  the KME continues to produce keys. Sezar should observe
  \texttt{psk\_age\_seconds} rising past policy.
\item
  \textbf{R5 --- Bifurcated SAE.} The strongSwan endpoint sees a healthy
  KME while the Wireguard endpoint sees a failed KME (simulating partial
  KME outage). Sezar should report inconsistent per-session
  \texttt{channel\_protection} while the \texttt{qkd\_link} aggregate is
  \texttt{degraded}.
\end{itemize}

\subsubsection{8.2.2 Metrics}\label{metrics-1}

We measure: SAE failover correctness, Sezar observation latency
(emulator change → emitted event), event-ordering correctness under
concurrent KME polls, and posture-rollup correctness (particularly that
the unified \(q\) score reflects the degradation appropriately).

\subsubsection{8.2.3 Results}\label{results}

All five scenarios ran end-to-end in compressed time (30--60s each
rather than the full 4--24h documented duration; the runner is at
\texttt{studies/study2/run.sh}). Per-scenario captures and the analysis
script are at \texttt{studies/study2/}.

\textbf{Classification correctness: 13/13.} Every operator-induced state
change produced a downstream \texttt{link\_health} reading matching the
expected post-op state.

\begin{longtable}[]{@{}lrrr@{}}
\toprule\noalign{}
Scenario & Events captured & Induced transitions & Matched \\
\midrule\noalign{}
\endhead
\bottomrule\noalign{}
\endlastfoot
R1 --- steady-state & 33 & 1 & 1/1 \\
R2 --- ramp & 63 & 5 & 5/5 \\
R3 --- hard-failure & 48 & 3 & 3/3 \\
R4 --- stale-PSK & 33 & 1 & 1/1 \\
R5 --- bifurcated & 48 & 3 & 3/3 \\
\end{longtable}

\textbf{Observation latency: p50 = 0.71s, range 0.70--0.71s} across all
13 transitions, against a configured 1-second poll interval. The latency
band hovers tightly around half the poll period --- the analytical
expectation for periodic polling --- and dominates over event-emission
cost. Increasing the poll interval moves the median latency linearly; in
real ETSI 014 deployments where the operational poll cadence is
typically 5--10 s {[}41{]}, this gives a p50 closer to 2.5--5 s.

The R3 link\_health timeline (plotted in the magazine companion;
analogue plots for every scenario ship in
\texttt{studies/study2/plots/}) shows the hard-failure transition
landing as expected.

R5 illustrates the per-session attribution case. When one KME returns
503, the link-level event flips to \texttt{failed} even though a
separately healthy paired KME continues to deliver keys. KME-only
telemetry cannot distinguish this from a full outage; the
channel-protection block on per-session events does.

R4 surfaces a gap we left in by design. Sezar's KME-side polling alone
cannot distinguish a fresh PSK from one the SAE has held for hours; the
\texttt{psk\_age\_seconds} field on \texttt{channel\_protection} is
populated only when the SAE co-operates. Closed-source SAEs without that
instrumentation are visible to Sezar only at the link layer. We
open-source patches for strongSwan, Wireguard, and a sample TLS endpoint
to populate the field.

The emulator, replay scripts, and analysis notebooks ship under MIT
alongside Sezar.

\subsection{8.3 Study 3 --- Axis G on Fifty Open-Source Server
Projects}\label{study-3-axis-g-on-fifty-open-source-server-projects}

\subsubsection{8.3.1 Methodology}\label{methodology-2}

We select fifty widely deployed open-source server projects spanning
HTTP, mail, database, message broker, and VPN categories: nginx, Apache
httpd, HAProxy, Envoy, Caddy, Traefik, OpenSSH server, Postfix, Dovecot,
Exim, PostgreSQL, MySQL, Redis, MongoDB, RabbitMQ, Kafka, OpenVPN,
strongSwan, Wireguard, PowerDNS, BIND, Unbound, CoreDNS, and others. The
full list, with hand-graded ground-truth levels and reviewer notes,
ships as \texttt{crates/sezar-agility/corpus/oss-50-v1.csv} in the Sezar
repository.

For each project we:

\begin{enumerate}
\def\labelenumi{\arabic{enumi}.}
\tightlist
\item
  Run \texttt{sezar-agility} against the source repository at a pinned
  commit.
\item
  Run \texttt{sezar-agility} against the installed binary on Rocky Linux
  10 with default package configuration.
\item
  Hand-grade the project against the §5.3 rubric, using two reviewers
  and reporting inter-rater agreement.
\item
  Report the divergence between automatic and hand-graded scores; treat
  the hand grade as ground truth for the corpus.
\end{enumerate}

\subsubsection{8.3.2 Metrics}\label{metrics-2}

We report per-project agility level, the per-evidence contribution to
the score, the false-negative and false-positive rates of the Semgrep
ruleset against the ground-truth grading, and the distribution of
agility levels across the corpus by category.

\subsubsection{8.3.3 Results (n = 11 pilot
subset)}\label{results-n-11-pilot-subset}

The full OSS-50 corpus is committed at
\texttt{crates/sezar-agility/corpus/oss-50-v1.csv} with hand-graded
ground truth. A pilot run over an 11-project subset spanning nine
categories (HTTP, mail, DB, message-broker, DNS, VPN/secure-shell,
messaging, certificate-authority, time) exercises the full pipeline:

\begin{longtable}[]{@{}
  >{\raggedright\arraybackslash}p{(\linewidth - 8\tabcolsep) * \real{0.2192}}
  >{\raggedright\arraybackslash}p{(\linewidth - 8\tabcolsep) * \real{0.3151}}
  >{\raggedright\arraybackslash}p{(\linewidth - 8\tabcolsep) * \real{0.1918}}
  >{\raggedright\arraybackslash}p{(\linewidth - 8\tabcolsep) * \real{0.1918}}
  >{\raggedright\arraybackslash}p{(\linewidth - 8\tabcolsep) * \real{0.0822}}@{}}
\toprule\noalign{}
\begin{minipage}[b]{\linewidth}\raggedright
Project
\end{minipage} & \begin{minipage}[b]{\linewidth}\raggedright
Category
\end{minipage} & \begin{minipage}[b]{\linewidth}\raggedright
Hand-grade
\end{minipage} & \begin{minipage}[b]{\linewidth}\raggedright
Scanner
\end{minipage} & \begin{minipage}[b]{\linewidth}\raggedright
Match
\end{minipage} \\
\midrule\noalign{}
\endhead
\bottomrule\noalign{}
\endlastfoot
nginx & http\_server & configurable & configurable & yes \\
haproxy & http\_server & configurable & configurable & yes \\
caddy & http\_server & configurable & configurable & yes \\
unbound & dns\_server & configurable & configurable & yes \\
coredns & dns\_server & configurable & configurable & yes \\
postfix & mail\_server & configurable & configurable & yes \\
redis & database & configurable & configurable & yes \\
nats-server & message\_broker & configurable & configurable & yes \\
step-ca & certificate\_authority & configurable & configurable & yes \\
prosody & messaging & configurable & configurable & yes \\
wireguard-tools & vpn\_secure\_shell & pinned & pinned & yes \\
chrony & time & configurable & \textbf{pinned} & no \\
\end{longtable}

\textbf{Agreement: 10/11 (91\%); Cohen's \(\kappa = 0.62\)} (substantial
agreement on the Landis--Koch scale). The single dissent is
\textbf{chrony}, where the static evidence is genuinely ambiguous
between a \texttt{configurable} reading (NTS key-types are
operator-controlled in \texttt{chrony.conf}) and a \texttt{pinned}
reading (the NTP authentication path embeds many literal symmetric-
algorithm references). The scanner picked the latter on the strength of
11 hard-coded references; no capability rule fired on chrony's NTS
surface because it uses neither OpenSSL nor Go's \texttt{crypto/tls}. We
flag this for v2 of the ruleset and the corpus.

The confusion matrix (plotted in the magazine companion) visualises this
result. Per-project event JSON, evidence listings, and the full
agreement TSV are at \texttt{studies/study3/results/}.

Failure modes characterised during the run:

\begin{itemize}
\tightlist
\item
  \textbf{Go ecosystem coverage gap.} The v1 ruleset's first cut matched
  only OpenSSL-style API calls (\texttt{SSL\_CTX\_set\_cipher\_list},
  \texttt{SSL\_CONF\_cmd}, \ldots). Three projects (coredns, step-ca,
  caddy) initially classified as \texttt{pinned} because nothing
  matched. Adding a one-rule \texttt{go-stdlib-tls-config} pattern
  matching \texttt{crypto/tls}'s \texttt{Config} / \texttt{CipherSuites}
  / \texttt{MinVersion} fields recovered all three to
  \texttt{configurable}.
\item
  \textbf{Aggregation policy.} Conservative-min aggregation
  systematically misclassified nginx: nginx embeds literal algorithm
  names in source (for crypto-impl code paths and tests, emit\_level
  \texttt{pinned}) and exposes operator config knobs (emit\_level
  \texttt{configurable}). The §7.3.3 most-agile- wins policy resolves
  the ambiguity in favour of the capability surface; conservative-min
  remains available as a regulatory-context alternative.
\item
  \textbf{No-evidence fallback.} Wireguard correctly classified
  \texttt{pinned} with zero evidence collected --- the project's
  userspace tooling has no OpenSSL or Go-tls surface, so the documented
  \texttt{UNKNOWN\_LEVEL\_FALLBACK\ =\ Pinned} applies. The ground-truth
  grade agrees.
\end{itemize}

Scaling the pilot to the full OSS-50 list is mechanical; the runner
already iterates over the corpus CSV.

\subsection{8.4 End-to-end pipeline}\label{end-to-end-pipeline}

\texttt{scripts/demo.sh} exercises the full V1 pipeline end-to-end: boot
the KME emulator, the QKD collector, sezar-server, and seed events from
both the bundled zgrab2 fixture and a synthetic FIPS-locked asset. The
collector's \texttt{/v1/posture} endpoint returns:

\begin{Shaded}
\begin{Highlighting}[]
\FunctionTok{\{}
  \DataTypeTok{"org\_q"}\FunctionTok{:} \FloatTok{0.627}\FunctionTok{,}
  \DataTypeTok{"deadline"}\FunctionTok{:} \StringTok{"2030{-}01{-}01T00:00:00Z"}\FunctionTok{,}
  \DataTypeTok{"horizon\_years"}\FunctionTok{:} \FloatTok{5.0}\FunctionTok{,}
  \DataTypeTok{"assets"}\FunctionTok{:} \DecValTok{5}\FunctionTok{,}
  \DataTypeTok{"blocked\_count"}\FunctionTok{:} \DecValTok{1}
\FunctionTok{\}}
\end{Highlighting}
\end{Shaded}

The inventory ordering, sorted by \(q\) descending, places the TLS 1.0 +
RC4 + SHA-1 legacy host at the top (q ≈ 0.72), the TLS 1.2 ECDHE+RSA
classical host just below (q ≈ 0.69), the FIPS-locked appliance third
with the \texttt{BLOCKED} flag raised (q ≈ 0.68), the modern TLS 1.3 +
ECDSA-P256 hybrid-PSK QKD host at the bottom of the priority queue (q ≈
0.43). This matches the model's expected behaviour from §5.4 (axis
combination) and §5.5 (edge cases). The channel-protection axis pulls
the QKD-protected asset down even though its algorithmic content is
classical; the agility axis pushes the FIPS-locked asset up even though
its observed primitives match the nominally-agile modern host; the
deadline-tension term leaves the legacy host at the top of the queue.

The implementation's pure-Rust rollup (in
\texttt{crates/sezar-server/src/posture.rs}) was verified against the
magazine companion's four-asset worked example: the unit tests
\texttt{worked\_example\_alpha\_q\_matches\_paper} and
\texttt{worked\_example\_delta\_q\_matches\_paper} assert that the
implementation reproduces α = 0.544 (modern-agile host, no QKD,
evaluated 2026-05-13) and δ = 0.392 (legacy classical host, same date)
to within 0.01.

To our knowledge no comparable open pipeline covers all three axes
end-to-end on commodity hardware. The Tranco-1k scan reported in §8.1 is
the broad-corpus result for Axis A; scaling Study 3 from the 11-project
pilot to the full OSS-50 corpus is the remaining mechanical step on the
published runners.

\begin{center}\rule{0.5\linewidth}{0.5pt}\end{center}

\section{9. Discussion}\label{discussion}

\subsection{9.1 Limitations}\label{limitations}

Our agility scoring is static and pattern-based; it is necessarily
approximate. A project may \emph{appear} pinned in source while actually
being agile through a runtime extension mechanism we did not
pattern-match. Conversely, a project may appear agile via a config field
that is in practice never changed. We address false negatives by
reporting per-evidence detail; we cannot fully address the second class
without operator input.

Our channel-protection axis depends on cooperating SAE instrumentation
for per-session attribution. Sezar can observe the KME state
independently, but linking a specific TLS or IPsec session to a specific
consumed key requires the SAE to emit the \texttt{key\_id\_observed}
field. We expect adoption in cooperating software (we open-source
patches for strongSwan, Wireguard, and a sample TLS endpoint as part of
the release) but note that closed-source SAEs may report only at the
link level.

Our threat model accepts the NIST PQC hardness assumptions and the
standard QKD security argument. Compromise of either --- by mathematical
break (PQC) or by implementation attack (QKD) --- would require
recalibration of the scoring tables. The schema supports this: the
rollup constants are operator-tunable in configuration, not compiled in.

We do not address economic or political constraints on migration (budget
cycles, regulatory approval lag, vendor support windows). These are
first-order operator concerns and we acknowledge them as out of scope
for the observability layer; they are downstream consumers of the
posture data.

\subsection{9.2 Ethical Considerations}\label{ethical-considerations}

Active TLS scanning, even of public web hosts, must be conducted
responsibly. We follow established practice from prior measurement work
{[}38{]}: one connection per host, clear User-Agent identification,
opt-out instructions linked from the identifying URL, conservative scan
rate, no protocol downgrade, no repeated probing. The Tranco list
excludes adult content and known sensitive categories; we additionally
exclude any host that returns a robots.txt directive on its base URL
within the first 1024 bytes.

The agility scanner operates on source code already public under
open-source licenses and on installed binaries on hosts the operator
controls. No private code or vendor materials are processed.

The QKD emulator generates synthetic key material only; it is not
connected to a real QKD link in our published study. Operators
integrating real KMEs are responsible for the security of those KMEs and
the surrounding network, including mutual TLS authentication of the SAE.

\subsection{9.3 Deployment Guidance}\label{deployment-guidance}

For operators planning to deploy Sezar against a real environment, we
recommend a phased rollout:

\begin{enumerate}
\def\labelenumi{\arabic{enumi}.}
\tightlist
\item
  \textbf{Phase 1: Inventory only.} Run \texttt{sezar-net} and
  \texttt{sezar-agility} in observe-only mode. Establish a baseline.
\item
  \textbf{Phase 2: Add deadline.} Configure \(D\) per applicable
  regulatory regime. Observe how \(q\) evolves with no other change.
\item
  \textbf{Phase 3: Prioritize.} Sort assets by \(q\) descending, migrate
  the top \(N\) by quarter.
\item
  \textbf{Phase 4: Integrate QKD telemetry where it exists.} Add
  \texttt{sezar-qkd} only when a real KME is present and the SAE
  instrumentation is in place; the channel-protection axis is
  \emph{additive} and never required.
\end{enumerate}

We caution against treating \(q\) as a primary KPI. It is a
\emph{comparative} score, calibrated for relative prioritization within
an environment. Inter-organization comparison requires shared \(D\),
shared weights, and shared corpora.

\subsection{9.4 Future Work}\label{future-work}

Several extensions are natural and intentionally deferred:

\begin{itemize}
\tightlist
\item
  \textbf{Time-decay.} Treating an observation made three months ago as
  less authoritative than one made today.
\item
  \textbf{Asset relationship modeling.} Linking a TLS session's
  certificate to the issuing CA's signing key allows the posture of an
  upstream asset to propagate into the downstream score.
\item
  \textbf{Adversary modeling.} Allowing operators to plug in their own
  estimate of CRQC arrival as a probability distribution over \(D\)
  rather than a single date.
\item
  \textbf{Standardization.} Submitting the
  \texttt{crypto\_inventory\_event} schema to IETF as an Informational
  draft so other tools can emit and consume it.
\end{itemize}

\begin{center}\rule{0.5\linewidth}{0.5pt}\end{center}

\section{10. Conclusion}\label{conclusion}

Operators preparing for the post-quantum transition cannot plan from a
single bit per asset. The ``is my asset PQ-ready?'' framing collapses
two assets with very different remediation costs into one bucket and
ignores the channel through which their key material travels. We have
set out a three-axis posture model, an open event schema, an
implementation, and three empirical studies that anchor the baseline.
The combined score \(q\) exists to surface what the single-bit view
hides; the \texttt{BLOCKED} flag exists to catch what \(q\) alone does
not.

We do not claim the scoring constants are optimal, that the agility
rubric covers every deployment pattern, or that our threat model
survives every adversary. We do claim that surfacing all three axes
beats single-axis dashboards, and that the released artifacts let the
community critique, refine, and supersede our specific choices on shared
empirical footing.

\begin{center}\rule{0.5\linewidth}{0.5pt}\end{center}

\section{Acknowledgments}\label{acknowledgments}

The authors thank the operators and reviewers who tested early versions
of the Sezar runner against their own infrastructures and whose feedback
shaped the published ruleset and the agility rubric. The Tranco list
project provides the snapshot identifiers that make our broad-corpus
measurements reproducible. Any errors are our own.

\begin{center}\rule{0.5\linewidth}{0.5pt}\end{center}

\section{Author Contributions}\label{author-contributions}

\textbf{Aleaddin Özer:} Conceptualization, methodology, software,
investigation, data curation, writing --- original draft, writing ---
review and editing, project administration. \textbf{Murat Aydos:}
Conceptualization, methodology, validation, formal analysis, writing ---
review and editing, supervision. Both authors read and approved the
final manuscript.

\begin{center}\rule{0.5\linewidth}{0.5pt}\end{center}

\section{Competing Interests}\label{competing-interests}

The authors declare no competing financial or non-financial interests.
E2E Solutions develops the Sezar reference implementation released under
the MIT License but derives no direct commercial revenue from its
publication. The authors have no financial stake in any of the
cryptographic stacks, QKD vendors, or scanning tools discussed in this
work.

\begin{center}\rule{0.5\linewidth}{0.5pt}\end{center}

\section{Funding}\label{funding}

This work was self-funded by E2E Solutions and Hacettepe University. No
external grants, sponsorships, or contracts supported the design,
execution, or reporting of the research presented here.

\begin{center}\rule{0.5\linewidth}{0.5pt}\end{center}

\section{Data and Code Availability}\label{data-and-code-availability}

All artifacts described in this paper are released under the MIT License
at \url{https://github.com/e2esolutions-tech/sezar}. The repository
contains the \texttt{crypto\_inventory\_event\ v1.1} schema, the Sezar
reference implementation (network, certificate, blockchain,
key-management, QKD agents plus collector and dashboard), the ETSI GS
QKD 014 emulator and its replay corpus, the Semgrep crypto-agility
ruleset, the hand-graded ground-truth corpus of fifty open-source server
projects (\texttt{crates/sezar-agility/corpus/oss-50-v1.csv}), and the
runner scripts that reproduce Studies 1--3. Tranco snapshot
\texttt{6G8PX} (2026-05-13) was used for the Study 1 sample frame. Raw
NDJSON captures and analysis plots for Studies 1 and 2 are committed
under \texttt{studies/study1/} and \texttt{studies/study2/}; Study 3's
per-project evidence listings and the agreement TSV are under
\texttt{studies/study3/results/}. No human subjects, personally
identifying information, or proprietary data are involved.

\begin{center}\rule{0.5\linewidth}{0.5pt}\end{center}

\section{Corresponding Author}\label{corresponding-author}

Aleaddin Özer
(\href{mailto:info@e2esolutions.tech}{\nolinkurl{info@e2esolutions.tech}}).

\begin{center}\rule{0.5\linewidth}{0.5pt}\end{center}

\section{Author Bios}\label{author-bios}

\textbf{Aleaddin Özer} is Chief Information Officer at E2E Solutions,
where he leads enterprise cryptographic infrastructure and post-quantum
migration programs across critical-sector clients. His operational work
informs the Sezar reference implementation. Contact:
\href{mailto:info@e2esolutions.tech}{\nolinkurl{info@e2esolutions.tech}}.
ORCID:
\href{https://orcid.org/0000-0001-9389-5357}{0000-0001-9389-5357}.

\textbf{Murat Aydos} is Associate Professor at Hacettepe University,
working on applied cryptography, network security, and post-quantum
migration strategy. His research has informed this paper's threat-model
framing and crypto-agility scoring rubric. ORCID:
\href{https://orcid.org/0000-0002-7570-9204}{0000-0002-7570-9204}.

\protect\phantomsection\label{refs}
\begin{CSLReferences}{0}{0}
\bibitem[\citeproctext]{ref-shor1997}
\CSLLeftMargin{{[}1{]} }%
\CSLRightInline{P. W. Shor, {``Polynomial-time algorithms for prime
factorization and discrete logarithms on a quantum computer,''}
\emph{SIAM Journal on Computing}, vol. 26, no. 5, pp. 1484--1509, 1997,
doi:
\href{https://doi.org/10.1137/S0097539795293172}{10.1137/S0097539795293172}.}

\bibitem[\citeproctext]{ref-grover1996}
\CSLLeftMargin{{[}2{]} }%
\CSLRightInline{L. K. Grover, {``A fast quantum mechanical algorithm for
database search,''} in \emph{Proceedings of the 28th annual ACM
symposium on theory of computing (STOC)}, 1996, pp. 212--219. doi:
\href{https://doi.org/10.1145/237814.237866}{10.1145/237814.237866}.}

\bibitem[\citeproctext]{ref-nsa-cnsa2}
\CSLLeftMargin{{[}3{]} }%
\CSLRightInline{NSA, {``{Announcing the Commercial National Security
Algorithm Suite 2.0},''} U.S. National Security Agency, Sep. 2022.}

\bibitem[\citeproctext]{ref-nist-ir8547}
\CSLLeftMargin{{[}4{]} }%
\CSLRightInline{NIST, {``{Transition to Post-Quantum Cryptography
Standards},''} National Institute of Standards; Technology, NIST IR 8547
(Initial Public Draft), Nov. 2024.}

\bibitem[\citeproctext]{ref-fips203}
\CSLLeftMargin{{[}5{]} }%
\CSLRightInline{NIST, {``{FIPS 203: Module-Lattice-Based
Key-Encapsulation Mechanism Standard (ML-KEM)},''} National Institute of
Standards; Technology, Aug. 2024.}

\bibitem[\citeproctext]{ref-fips204}
\CSLLeftMargin{{[}6{]} }%
\CSLRightInline{NIST, {``{FIPS 204: Module-Lattice-Based Digital
Signature Standard (ML-DSA)},''} National Institute of Standards;
Technology, Aug. 2024.}

\bibitem[\citeproctext]{ref-fips205}
\CSLLeftMargin{{[}7{]} }%
\CSLRightInline{NIST, {``{FIPS 205: Stateless Hash-Based Digital
Signature Standard (SLH-DSA)},''} National Institute of Standards;
Technology, Aug. 2024.}

\bibitem[\citeproctext]{ref-ncsc-pqc}
\CSLLeftMargin{{[}8{]} }%
\CSLRightInline{UK NCSC, {``{Preparing for Quantum-Safe
Cryptography}.''} UK National Cyber Security Centre guidance, 2023.}

\bibitem[\citeproctext]{ref-anssi-pqc}
\CSLLeftMargin{{[}9{]} }%
\CSLRightInline{ANSSI, {``{ANSSI Views on the Post-Quantum Cryptography
Transition},''} Agence nationale de la sécurité des systèmes
d'information (France), 2022.}

\bibitem[\citeproctext]{ref-cisco-pqc-discovery}
\CSLLeftMargin{{[}10{]} }%
\CSLRightInline{Cisco Systems, {``{Post-Quantum Cryptography at Cisco:
Trust Center Overview and Quantum-Safe WAN Whitepaper}.''}
\url{https://www.cisco.com/site/us/en/about/trust-center/post-quantum-cryptography.html},
2024.}

\bibitem[\citeproctext]{ref-ibm-pqc-discovery}
\CSLLeftMargin{{[}11{]} }%
\CSLRightInline{IBM, {``{IBM Guardium Quantum Safe and IBM Quantum Safe
Explorer}.''} \url{https://www.ibm.com/products/guardium-quantum-safe},
2024.}

\bibitem[\citeproctext]{ref-microsoft-cryptotracker}
\CSLLeftMargin{{[}12{]} }%
\CSLRightInline{A. Shemesh and J. Rutzer, {``{Building your
cryptographic inventory: A customer strategy for cryptographic posture
management}.''} Microsoft Security Blog, 2026-04-16,
\url{https://www.microsoft.com/en-us/security/blog/2026/04/16/building-your-cryptographic-inventory-a-customer-strategy-for-cryptographic-posture-management/},
Apr. 2026.}

\bibitem[\citeproctext]{ref-cloudflare-pq-deploy}
\CSLLeftMargin{{[}13{]} }%
\CSLRightInline{B. Westerbaan, {``{The state of the post-quantum
Internet}.''} Cloudflare Blog, 2024-03-05,
\url{https://blog.cloudflare.com/pq-2024/}, Mar. 2024.}

\bibitem[\citeproctext]{ref-bsi-migration}
\CSLLeftMargin{{[}14{]} }%
\CSLRightInline{BSI, {``{Migration to Post-Quantum Cryptography:
Recommendations for Action},''} Bundesamt für Sicherheit in der
Informationstechnik (Germany), 2023.}

\bibitem[\citeproctext]{ref-bb84}
\CSLLeftMargin{{[}15{]} }%
\CSLRightInline{C. H. Bennett and G. Brassard, {``Quantum cryptography:
Public key distribution and coin tossing,''} in \emph{Proceedings of the
IEEE international conference on computers, systems and signal
processing}, Bangalore, India, 1984, pp. 175--179.}

\bibitem[\citeproctext]{ref-secoqc}
\CSLLeftMargin{{[}16{]} }%
\CSLRightInline{M. Peev, C. Pacher, R. Alléaume, C. Barreiro, J. Bouda,
\emph{et al.}, {``The {SECOQC} quantum key distribution network in
vienna,''} \emph{New Journal of Physics}, vol. 11, no. 7, p. 075001,
2009.}

\bibitem[\citeproctext]{ref-tokyo-qkd}
\CSLLeftMargin{{[}17{]} }%
\CSLRightInline{M. Sasaki, M. Fujiwara, H. Ishizuka, \emph{et al.},
{``Field test of quantum key distribution in the {Tokyo} {QKD}
{Network},''} \emph{Optics Express}, vol. 19, no. 11, pp. 10387--10409,
2011.}

\bibitem[\citeproctext]{ref-euroqci}
\CSLLeftMargin{{[}18{]} }%
\CSLRightInline{European Commission, {``{European Quantum Communication
Infrastructure (EuroQCI) Initiative}.''}
\url{https://digital-strategy.ec.europa.eu/en/policies/european-quantum-communication-infrastructure-euroqci},
2023.}

\bibitem[\citeproctext]{ref-etsi-qkd-014}
\CSLLeftMargin{{[}19{]} }%
\CSLRightInline{ETSI ISG-QKD, {``{GS QKD 014 V1.1.1 --- Quantum Key
Distribution; Protocol and Data Format of REST-Based Key Delivery
API},''} ETSI, Feb. 2019.}

\bibitem[\citeproctext]{ref-idq-deployments}
\CSLLeftMargin{{[}20{]} }%
\CSLRightInline{N. Aquina \emph{et al.}, {``{A Critical Analysis of
Deployed Use Cases for Quantum Key Distribution and Comparison with
Post-Quantum Cryptography},''} \emph{EPJ Quantum Technology}, vol. 12,
no. 51, 2025, doi:
\href{https://doi.org/10.1140/epjqt/s40507-025-00350-5}{10.1140/epjqt/s40507-025-00350-5}.}

\bibitem[\citeproctext]{ref-toshiba-qkd}
\CSLLeftMargin{{[}21{]} }%
\CSLRightInline{M. Pittaluga \emph{et al.}, {``600-km repeater-like
quantum communications with dual-band stabilization,''} \emph{Nature
Photonics}, vol. 15, pp. 530--535, 2021, doi:
\href{https://doi.org/10.1038/s41566-021-00811-0}{10.1038/s41566-021-00811-0}.}

\bibitem[\citeproctext]{ref-qti-qkd}
\CSLLeftMargin{{[}22{]} }%
\CSLRightInline{Quantum Telecommunications Italy (QTI), {``{QUID:
Italian National Quantum Communication Infrastructure Deployment}.''}
EuroQCI Italy project, \url{https://quid-euroqci-italy.eu/the-project/},
2024.}

\bibitem[\citeproctext]{ref-nsa-qkd-position}
\CSLLeftMargin{{[}23{]} }%
\CSLRightInline{NSA, {``{Quantum Key Distribution (QKD) and Quantum
Cryptography (QC)}.''} NSA Information Sheet, 2020.}

\bibitem[\citeproctext]{ref-anssi-qkd}
\CSLLeftMargin{{[}24{]} }%
\CSLRightInline{ANSSI, {``{Should Quantum Key Distribution Be Used for
Secure Communications?}''} Agence nationale de la sécurité des systèmes
d'information (France), 2020.}

\bibitem[\citeproctext]{ref-rfc7696}
\CSLLeftMargin{{[}25{]} }%
\CSLRightInline{R. Housley, {``{Guidelines for Cryptographic Algorithm
Agility and Selecting Mandatory-to-Implement Algorithms},''} IETF, RFC
7696, 2015.}

\bibitem[\citeproctext]{ref-nist-pqc-call}
\CSLLeftMargin{{[}26{]} }%
\CSLRightInline{NIST, {``{Post-Quantum Cryptography Standardization Call
for Proposals}.''}
\url{https://csrc.nist.gov/projects/post-quantum-cryptography}, 2016.}

\bibitem[\citeproctext]{ref-nist-fn-dsa}
\CSLLeftMargin{{[}27{]} }%
\CSLRightInline{R. Perlner, {``{FIPS 206 Status Update: FN-DSA
(Falcon)},''} National Institute of Standards; Technology, Computer
Security Division, 2025.}

\bibitem[\citeproctext]{ref-nist-pqc-onramp}
\CSLLeftMargin{{[}28{]} }%
\CSLRightInline{NIST, {``{Additional Digital Signature Schemes and KEM
On-Ramp Round}.''} NIST Post-Quantum Cryptography project page, 2023.}

\bibitem[\citeproctext]{ref-chrome-pq-tls}
\CSLLeftMargin{{[}29{]} }%
\CSLRightInline{D. Adrian, B. Beck, D. Benjamin, and D. O'Brien,
{``{Advancing Our Amazing Bet on Asymmetric Cryptography}.''} Chromium
Blog, 2024-05-23 (Chrome 131 shipped ML-KEM Nov 2024),
\url{https://blog.google/chromium/advancing-our-amazing-bet-on-asymmetric/},
May 2024.}

\bibitem[\citeproctext]{ref-firefox-pq-tls}
\CSLLeftMargin{{[}30{]} }%
\CSLRightInline{Mozilla, {``{Firefox 132 Release Notes: Post-Quantum Key
Exchange (mlkem768x25519) for TLS 1.3}.''}
\url{https://www.firefox.com/en-US/firefox/132.0/releasenotes/}, Oct.
2024.}

\bibitem[\citeproctext]{ref-pq-tls-measurement}
\CSLLeftMargin{{[}31{]} }%
\CSLRightInline{Cloudflare Research, {``{State of the Post-Quantum
Internet in 2025}.''} Cloudflare Blog,
\url{https://blog.cloudflare.com/pq-2025/}, 2025.}

\bibitem[\citeproctext]{ref-digicert-pq}
\CSLLeftMargin{{[}32{]} }%
\CSLRightInline{DigiCert, {``{ML-DSA and FN-DSA Post-Quantum Signature
Readiness}.''} DigiCert PQC Insights and Blog,
\url{https://www.digicert.com/insights/post-quantum-cryptography/mldsa},
2025.}

\bibitem[\citeproctext]{ref-micius}
\CSLLeftMargin{{[}33{]} }%
\CSLRightInline{S.-K. Liao, W.-Q. Cai, W.-Y. Liu, \emph{et al.},
{``Satellite-to-ground quantum key distribution,''} \emph{Nature}, vol.
549, pp. 43--47, 2017.}

\bibitem[\citeproctext]{ref-twin-field}
\CSLLeftMargin{{[}34{]} }%
\CSLRightInline{M. Lucamarini, Z. L. Yuan, J. F. Dynes, and A. J.
Shields, {``Overcoming the rate--distance limit of quantum key
distribution without quantum repeaters,''} \emph{Nature}, vol. 557, pp.
400--403, 2018.}

\bibitem[\citeproctext]{ref-hybrid-tls-psk}
\CSLLeftMargin{{[}35{]} }%
\CSLRightInline{T. Wiggers, S. Celi, P. Schwabe, D. Stebila, and N.
Sullivan, {``{KEM-based pre-shared-key handshakes for TLS 1.3}.''} IETF
Internet-Draft draft-wiggers-tls-authkem-psk-05 (2026-05-04, individual
submission),
\url{https://datatracker.ietf.org/doc/draft-wiggers-tls-authkem-psk/},
May 2026.}

\bibitem[\citeproctext]{ref-nist-sp-1800-38a}
\CSLLeftMargin{{[}36{]} }%
\CSLRightInline{NIST National Cybersecurity Center of Excellence,
{``{Migration to Post-Quantum Cryptography (NIST SP 1800-38, Vols.~A--C,
Preliminary Drafts)},''} National Institute of Standards; Technology,
2023.}

\bibitem[\citeproctext]{ref-fips140-3}
\CSLLeftMargin{{[}37{]} }%
\CSLRightInline{NIST, {``{FIPS 140-3: Security Requirements for
Cryptographic Modules},''} National Institute of Standards; Technology,
2019.}

\bibitem[\citeproctext]{ref-durumeric-tls}
\CSLLeftMargin{{[}38{]} }%
\CSLRightInline{Z. Durumeric, J. Kasten, M. Bailey, and J. A. Halderman,
{``Analysis of the {HTTPS} certificate ecosystem,''} in
\emph{Proceedings of the ACM internet measurement conference (IMC)},
2013.}

\bibitem[\citeproctext]{ref-holz-tls}
\CSLLeftMargin{{[}39{]} }%
\CSLRightInline{R. Holz, L. Braun, N. Kammenhuber, and G. Carle, {``The
{SSL} landscape: A thorough analysis of the {X.509} {PKI} using active
and passive measurements,''} in \emph{Proceedings of the ACM internet
measurement conference (IMC)}, 2011.}

\bibitem[\citeproctext]{ref-felt-tls}
\CSLLeftMargin{{[}40{]} }%
\CSLRightInline{A. P. Felt, R. Barnes, A. King, C. Palmer, C. Bentzel,
and P. Tabriz, {``Measuring {HTTPS} adoption on the web,''} in
\emph{Proceedings of the 26th USENIX security symposium}, 2017.}

\bibitem[\citeproctext]{ref-etsi-qkd-002}
\CSLLeftMargin{{[}41{]} }%
\CSLRightInline{ETSI ISG-QKD, {``{GS QKD 002 --- Use Cases},''} ETSI,
2010.}

\bibitem[\citeproctext]{ref-etsi-qkd-008}
\CSLLeftMargin{{[}42{]} }%
\CSLRightInline{ETSI ISG-QKD, {``{GS QKD 008 --- Module Security
Specification},''} ETSI, 2010.}

\bibitem[\citeproctext]{ref-etsi-qkd-011}
\CSLLeftMargin{{[}43{]} }%
\CSLRightInline{ETSI ISG-QKD, {``{GS QKD 011 --- Component
Characterization: Characterizing Optical Components for QKD Systems},''}
ETSI, 2016.}

\bibitem[\citeproctext]{ref-ncc-group-crypto-audits}
\CSLLeftMargin{{[}44{]} }%
\CSLRightInline{NCC Group, {``{Cryptographic Implementation Audits and
Crypto-Agility Findings (Public Reports)}.''} NCC Group Research blog,
\url{https://research.nccgroup.com/}, 2018-\/-2024.}

\bibitem[\citeproctext]{ref-owasp-crypto}
\CSLLeftMargin{{[}45{]} }%
\CSLRightInline{OWASP Foundation, {``{OWASP Cryptographic Storage Cheat
Sheet}.''}
\url{https://cheatsheetseries.owasp.org/cheatsheets/Cryptographic_Storage_Cheat_Sheet.html},
2024.}

\bibitem[\citeproctext]{ref-tranco}
\CSLLeftMargin{{[}46{]} }%
\CSLRightInline{V. Le Pochat, T. Van Goethem, S. Tajalizadehkhoob, M.
Korczyński, and W. Joosen, {``{Tranco}: A research-oriented top sites
ranking hardened against manipulation,''} in \emph{Proceedings of the
network and distributed system security symposium (NDSS)}, 2019.}

\bibitem[\citeproctext]{ref-zgrab2}
\CSLLeftMargin{{[}47{]} }%
\CSLRightInline{ZMap Project, {``{zgrab2: A Fast Application-Layer
Network Scanner}.''} \url{https://github.com/zmap/zgrab2}, 2024.}

\end{CSLReferences}
